\section{Context Loop Wars: puntos de entrada, herramientas y retorno de datos}

Cuando un coding agent, un agent de oficina o un agent sectorial entra en un flujo de trabajo real, la empresa que controla el punto de entrada del usuario puede observar la distribución de las tareas, la empresa dueña del modelo puede controlar la capacidad y el precio, y la empresa dueña de las herramientas y del entorno define los límites de acción. Estas tres posiciones no siempre las ocupa la misma empresa, por lo que las relaciones de colaboración se convierten rápidamente en competencia. Esta sección trata los casos de Windsurf, OpenAI, Anthropic y Mistral que eligió el expositor como un análisis de estructura de sistemas, no como chisme empresarial.

\subsection{Qué reveló una disputa por una adquisición}

El expositor usa los cambios del mercado de coding agents en 2025 para mostrar el valor del context loop. Windsurf se ubica en el punto de entrada del flujo de trabajo de los desarrolladores, los proveedores de modelos aportan la capacidad de base, y el código de los usuarios junto con la retroalimentación de la ejecución forman un entorno de alto valor. En cuanto una de las partes teme que la otra interiorice los datos de interacción, la distribución o la relación con el producto, la colaboración original vía API da paso a restricciones de acceso, negociaciones de adquisición y despliegue de modelos alternativos.

\begin{figure}[H]
\centering
\includegraphics[width=0.94\textwidth]{images/00-27-34-context-loop-wars.jpg}
\caption{Primera etapa de las Context Loop Wars: el conflicto en torno al punto de entrada del flujo de trabajo de Windsurf y al acceso a los modelos.}
\figsource{Video oficial de Stanford Online 00:27:34, `images/00-27-34-context-loop-wars.jpg`.}
\end{figure}

\begin{knowledgebox}{Lectura de la figura: no hay que fijarse solo en el precio de la adquisición}
Lo más importante son cuatro flujos: por dónde entran las solicitudes del usuario, dónde se ejecutan el código y las herramientas, qué empresa aporta el modelo que hace la inferencia y a qué sistema de datos regresan los resultados y las modificaciones. La empresa del punto de entrada controla la demanda y la distribución; la empresa del modelo controla la capacidad y la oferta; la empresa de nube o de herramientas controla la infraestructura de ejecución. Una transacción solo reorganiza el control de estos cuatro flujos. Una vez entendida la estructura, el método de análisis sigue siendo válido aunque cambien las empresas concretas.
\end{knowledgebox}

\begin{figure}[H]
\centering
\includegraphics[width=0.94\textwidth]{images/00-29-04-context-loop-wars-cont.jpg}
\caption{Continuación de las Context Loop Wars: modelos, nube, herramientas de desarrollo y talento se recombinan dentro del mismo ciclo cerrado.}
\figsource{Video oficial de Stanford Online 00:29:04, `images/00-29-04-context-loop-wars-cont.jpg`.}
\end{figure}

La segunda figura amplía una transacción aislada a una relación de ecosistema. Al leerla conviene comparar qué le falta a cada participante: a los laboratorios de modelos pueden faltarles una distribución estable y entornos de alta calidad; a las herramientas de desarrollo, modelos propios y control del costo de inferencia a largo plazo; las plataformas de nube pueden querer fijar la demanda de cómputo dentro de su propia infraestructura, y el talento y los equipos de investigación se desplazan según dónde estén el capital y los datos. La llamada context war es, en esencia, una disputa por un ciclo cerrado de mejora sostenible, no por un lote de registros de chat estáticos.

\begin{warningbox}{Nota de clase: los casos ilustran mecanismos, no garantizan que la cronología siga vigente}
El estado de las fusiones y adquisiciones, las políticas de API y las capacidades de los productos cambian. Estas notas conservan los casos históricos que el expositor usó en esa sesión para explicar el valor del context. El lector debe verificar los hechos más recientes y después usar el marco «punto de entrada—entorno—modelo—retroalimentación—distribución» para juzgar la nueva situación; las conexiones entre empresas de las diapositivas no deben tomarse como un mapa permanente de la industria.
\end{warningbox}

\subsection{Mistral: el context clave también puede ser un entorno soberano}

El context no pertenece solo a un producto SaaS. El idioma, el derecho, las compras del sector público, los requisitos de seguridad nacional, la residencia de los datos y los ecosistemas industriales locales también forman entornos difíciles de replicar de una región a otra. El expositor explica la motivación detrás de la creación de Mistral a partir de las necesidades europeas de soberanía y de pesos abiertos: cuando las capacidades clave de los modelos dependen por completo de proveedores externos, a las instituciones locales les resulta difícil controlar los datos, los costos, la capacidad de auditoría y su poder de negociación a largo plazo.

\begin{figure}[H]
\centering
\includegraphics[width=0.94\textwidth]{images/00-30-00-mistral-critical-contexts.jpg}
\caption{El caso de Mistral: los idiomas, la regulación, las instituciones públicas y las necesidades industriales de Europa constituyen critical contexts/environments.}
\figsource{Video oficial de Stanford Online 00:30:00, `images/00-30-00-mistral-critical-contexts.jpg`.}
\end{figure}

\begin{knowledgebox}{Lectura de la figura: el significado sistémico de la sovereign AI}
La sovereign AI no consiste solo en «entrenar un modelo nacional». Abarca, como mínimo, datos de idioma y cultura locales, licencias de modelo aceptables, residencia de los datos, control de la cadena de suministro, interfaces de despliegue para el sector público, talento y cómputo, así como la capacidad de mantener el servicio en una crisis. Los pesos abiertos pueden mejorar la posibilidad de inspección y las opciones de despliegue, pero no resuelven por sí solos el costo de entrenamiento, la responsabilidad en seguridad ni la fragmentación del ecosistema. Al final, las aspiraciones de soberanía tienen que concretarse en infraestructura operable y en arreglos institucionales.
\end{knowledgebox}

\begin{importantbox}{Las tres escaseces del context}
El context de producto es escaso porque los flujos de trabajo reales y la confianza de los usuarios son difíciles de replicar; el context organizacional es escaso porque el conocimiento privado, los permisos y las decisiones históricas no pueden recopilarse públicamente; el context soberano es escaso porque el idioma, el derecho y el interés público son territoriales. Cuanto más se parezcan las capacidades de los modelos, más probable será que quien pueda entrar de forma segura en estos entornos obtenga la siguiente ronda de retroalimentación de alto valor.
\end{importantbox}

\subsection{Resumen de la sección}

Las context loop wars concretan la abstracta «ventaja de datos» como el control sobre el punto de entrada, las herramientas, el entorno de ejecución, la oferta de modelos y el retorno de la retroalimentación. El caso de Windsurf muestra el ciclo cerrado de producto; el caso de Mistral, los entornos soberanos y públicos. La siguiente sección plantea otra pregunta: aun con el entorno y la retroalimentación, ¿puede el RL seguir escalando de forma sostenida, o se topará con la frontera entre lo medible y lo no medible?
